\chapter{导数及其应用}

\section{导数的概念}

\begin{definition}[导数]
设函数 $y = f(x)$ 在点 $x_0$ 的某个邻域内有定义，当自变量 $x$ 在 $x_0$ 处有增量 $\Delta x$ 时，函数有相应的增量 $\Delta y = f(x_0 + \Delta x) - f(x_0)$. 如果当 $\Delta x \to 0$ 时，比值 $\dfrac{\Delta y}{\Delta x}$ 的极限存在，则称函数 $f(x)$ 在点 $x_0$ 处\textbf{可导}，并称此极限为函数 $f(x)$ 在点 $x_0$ 处的\textbf{导数}，记作 $f'(x_0)$ 或 $\left.\dfrac{dy}{dx}\right|_{x=x_0}$：
\[ f'(x_0) = \lim_{\Delta x \to 0} \frac{f(x_0 + \Delta x) - f(x_0)}{\Delta x} \]
\end{definition}

\begin{note}
导数的几何意义：函数 $y = f(x)$ 在点 $x_0$ 处的导数 $f'(x_0)$ 就是曲线 $y = f(x)$ 在点 $(x_0, f(x_0))$ 处的切线斜率。
\end{note}

\begin{example}
用导数定义求函数 $f(x) = x^2$ 在 $x = 1$ 处的导数。
\end{example}

\textbf{解.}
\begin{align*}
f'(1) &= \lim_{\Delta x \to 0} \frac{(1 + \Delta x)^2 - 1^2}{\Delta x} \\
&= \lim_{\Delta x \to 0} \frac{1 + 2\Delta x + (\Delta x)^2 - 1}{\Delta x} \\
&= \lim_{\Delta x \to 0} (2 + \Delta x) = 2
\end{align*}

因此，$f'(1) = 2$. \hfill $\square$

\section{基本初等函数的导数公式}

\begin{theorem}[基本求导公式]
\begin{enumerate}
    \item $(C)' = 0$ \quad （$C$ 为常数）
    \item $(x^n)' = nx^{n-1}$ \quad （$n \in \mathbb{Q}$）
    \item $(\sin x)' = \cos x$
    \item $(\cos x)' = -\sin x$
    \item $(e^x)' = e^x$
    \item $(\ln x)' = \dfrac{1}{x}$ \quad （$x > 0$）
\end{enumerate}
\end{theorem}

\begin{exercise}
求下列函数的导数：
\begin{enumerate}
    \item $y = x^5$
    \item $y = \dfrac{1}{x^2}$
    \item $y = \sqrt{x}$
\end{enumerate}
\end{exercise}

\textbf{解析：}
\begin{enumerate}
    \item $y' = 5x^4$
    \item $y = x^{-2}$，故 $y' = -2x^{-3} = -\dfrac{2}{x^3}$
    \item $y = x^{1/2}$，故 $y' = \dfrac{1}{2}x^{-1/2} = \dfrac{1}{2\sqrt{x}}$
\end{enumerate}

\section{导数的运算法则}

\begin{theorem}[导数的四则运算法则]
设函数 $u(x)$ 和 $v(x)$ 都可导，则：
\begin{enumerate}
    \item $(u \pm v)' = u' \pm v'$
    \item $(uv)' = u'v + uv'$ \quad （乘法法则）
    \item $\left(\dfrac{u}{v}\right)' = \dfrac{u'v - uv'}{v^2}$ \quad （$v \neq 0$，除法法则）
\end{enumerate}
\end{theorem}

\begin{example}
求函数 $y = x^3 \cdot \ln x$ 的导数。
\end{example}

\textbf{解.} 令 $u = x^3$，$v = \ln x$，则 $u' = 3x^2$，$v' = \dfrac{1}{x}$.

由乘法法则：
\[ y' = 3x^2 \cdot \ln x + x^3 \cdot \frac{1}{x} = 3x^2\ln x + x^2 = x^2(3\ln x + 1) \]
\hfill $\square$

\section{利用导数研究函数的单调性}

\begin{theorem}[单调性判定]
设函数 $y = f(x)$ 在区间 $(a, b)$ 内可导：
\begin{enumerate}
    \item 若 $f'(x) > 0$ 对任意 $x \in (a, b)$ 成立，则 $f(x)$ 在 $(a, b)$ 上单调递增；
    \item 若 $f'(x) < 0$ 对任意 $x \in (a, b)$ 成立，则 $f(x)$ 在 $(a, b)$ 上单调递减。
\end{enumerate}
\end{theorem}

\begin{note}
记忆口诀：\textbf{正增负减}——导数为正则函数增，导数为负则函数减。
\end{note}

\begin{longexample}
求函数 $f(x) = x^3 - 3x^2$ 的单调区间。
\end{longexample}

\textbf{解.} 

\textbf{第一步：求导数}
\[ f'(x) = 3x^2 - 6x = 3x(x - 2) \]

\textbf{第二步：求临界点}
令 $f'(x) = 0$，得 $x = 0$ 或 $x = 2$.

\textbf{第三步：列表分析}

\begin{center}
\begin{tabular}{c|ccccc}
\toprule
$x$ & $(-\infty, 0)$ & $0$ & $(0, 2)$ & $2$ & $(2, +\infty)$ \\
\midrule
$f'(x)$ & $+$ & $0$ & $-$ & $0$ & $+$ \\
$f(x)$ & $\nearrow$ & 极大值 & $\searrow$ & 极小值 & $\nearrow$ \\
\bottomrule
\end{tabular}
\end{center}

\textbf{第四步：得出结论}

函数 $f(x)$ 在 $(-\infty, 0)$ 和 $(2, +\infty)$ 上单调递增，在 $(0, 2)$ 上单调递减。\hfill $\square$

\section{函数的极值与最值}

\begin{definition}[极值]
设函数 $f(x)$ 在点 $x_0$ 的某个邻域内有定义：
\begin{enumerate}
    \item 若对邻域内任意 $x \neq x_0$，都有 $f(x) < f(x_0)$，则称 $f(x_0)$ 为函数的\textbf{极大值}，$x_0$ 为\textbf{极大值点}；
    \item 若对邻域内任意 $x \neq x_0$，都有 $f(x) > f(x_0)$，则称 $f(x_0)$ 为函数的\textbf{极小值}，$x_0$ 为\textbf{极小值点}。
\end{enumerate}
\end{definition}

\begin{theorem}[极值的必要条件]
若函数 $f(x)$ 在点 $x_0$ 处可导且取得极值，则 $f'(x_0) = 0$.
\end{theorem}

\begin{note}
使 $f'(x) = 0$ 的点称为函数的\textbf{驻点}（或稳定点）。驻点不一定是极值点，需要进一步判断。
\end{note}

\begin{example}
求函数 $f(x) = x^3 - 3x$ 的极值。
\end{example}

\textbf{解.}

求导：$f'(x) = 3x^2 - 3 = 3(x+1)(x-1)$

令 $f'(x) = 0$，得 $x = -1$ 或 $x = 1$.

列表分析：
\begin{center}
\begin{tabular}{c|ccccc}
\toprule
$x$ & $(-\infty, -1)$ & $-1$ & $(-1, 1)$ & $1$ & $(1, +\infty)$ \\
\midrule
$f'(x)$ & $+$ & $0$ & $-$ & $0$ & $+$ \\
$f(x)$ & $\nearrow$ & 极大值 & $\searrow$ & 极小值 & $\nearrow$ \\
\bottomrule
\end{tabular}
\end{center}

计算：$f(-1) = (-1)^3 - 3(-1) = 2$，$f(1) = 1 - 3 = -2$.

因此，函数在 $x = -1$ 处取得极大值 $2$，在 $x = 1$ 处取得极小值 $-2$。\hfill $\square$
